% appendix_iv_validation.tex: Using an instrument observed on a subsample to
% validate a fixed-effects design (2026-09-03).
\section{An Instrument on a Subsample and the Design for the Full Sample}\label{appendix:ivvalidation}

This appendix states when an instrument that exists on a subsample $S$ of the data justifies a fixed-effects design on all of it, and how the standard errors should be computed. The argument is linear; the logit case follows by the same steps on the score.

\subsection{Setup}
For unit $i$ let $y_i = \tau x_i + c_i + e_i$, with $x_i$ the regressor of interest, $c_i$ an unobserved confounder correlated with $x_i$, and $e_i$ independent noise. A fixed-effects design $F$ partitions units into cells; $M_F = I - P_F$ is the within-cell projection. An instrument $z_i$ with $\mathbb{E}[z_i c_i \mid F] = \mathbb{E}[z_i e_i \mid F] = 0$ is observed only on $S$.

\subsection{What the fixed-effects estimator estimates}
On any sample $A$,
\begin{equation}
\hat\tau_F(A) = \tau + \underbrace{(x'M_F x)^{-1}x'M_F c}_{b_F(A)} + (x'M_F x)^{-1}x'M_F e,
\end{equation}
so the bias $b_F(A)$ is the within-cell regression coefficient of the confounder on $x$ in sample $A$. When cells nest within $S$ and $S^c$, the full-sample bias is the variance-weighted average
\begin{equation}\label{eq:biasavg}
b_F(\text{Full}) = w_S\, b_F(S) + (1 - w_S)\, b_F(S^c), \qquad w_S = \frac{x_S' M_F x_S}{x'M_F x}.
\end{equation}
The instrument identifies $\tau$ on $S$. A design selected because $\hat\tau_F(S) \approx \hat\tau_{IV}(S)$ and then applied to the full sample is unbiased if and only if $b_F(S^c) = 0$ given $b_F(S) = 0$. The difference estimator
\begin{equation}\label{eq:diffapp}
\hat\tau^{D} = \hat\tau_F(\text{Full}) - \big[\hat\tau_F(S) - \hat\tau_{IV}(S)\big]
\end{equation}
is unbiased if and only if $b_F(S^c) = b_F(S)$. The second condition contains the first as the case $b_F(S) = 0$; it does not require the design to be exactly right on $S$, and it uses the instrument as a level rather than as a test.

\subsection{The assumption}
Write $c = c_F + \eta$ with $c_F$ measurable with respect to the cells, so $M_F c = M_F \eta$. If the conditional law of $(x, \eta)$ given the cell is the same in $S$ and $S^c$, then $b_F(S) = b_F(S^c)$. This is unconfounded location in the sense of \citet{hotz_predicting_2005} and equi-confounding in the sense of \citet{li_identification_2024}: membership in $S$ is independent of the residual confounding mechanism given the cells, while cell means, effect sizes, and the distribution of $x$ across cells may differ. It is the weakest condition under which either procedure is consistent, and under it the difference estimator dominates selection: it needs no choice among designs and inherits no pre-test distortion \citep{leeb_model_2005, guggenberger_impact_2010}. A testable implication is that the estimated bias $\hat\tau_F(S) - \hat\tau_{IV}(S)$ is constant across groups of units within $S$; \citet{angrist_leveraging_2017} apply the same check to the bias of value-added estimates across schools with and without lotteries. Effects for groups that exist only in $S^c$ are identified only under the assumption itself or under a parametric bias function estimated on $S$ and extrapolated \citep{kallus_removing_2018}.

\subsection{Standard errors}
The three pieces of \eqref{eq:diffapp} share the observations in $S$, so
\begin{equation}
\mathrm{Var}(\hat\tau^{D}) = \mathrm{Var}\,\hat\tau_F(\text{Full}) + \mathrm{Var}\big[\hat\tau_F(S) - \hat\tau_{IV}(S)\big] - 2\,\mathrm{Cov}\big(\hat\tau_F(\text{Full}),\, \hat\tau_F(S) - \hat\tau_{IV}(S)\big),
\end{equation}
which a cluster bootstrap over the units of instrument variation, stratified by $S$, estimates directly; this is the variance of the prediction-powered estimator of \citet{angelopoulos_prediction_2023}. It exceeds the variance of the fixed-effects estimate alone by the variance of the estimated bias.


\subsection{The instruments, and why they fail}\label{appendix:failedinstruments}

For lender-years with at least fifty originations, 97 percent of loans, two leave-one-out instruments exist: the lender's share of long-term originations in the year, which shifts $T$ through the maturity menu \citep{argyle_monthly_2020, hertzberg_screening_2018}, and the lender's mean log loan size, which shifts $m$ given $T$ through the size menu. The two-stage least-squares payment coefficient on that subsample is negative, the opposite sign to every payment response in the paper, including the within-borrower estimate and the rule-based experiments (Table \ref{tab:identification}). Lender-years with larger loans serve safer borrowers, so the size-menu instrument moves the payment together with unobserved risk and the exclusion restriction fails; the bias the instrument would measure differs across regions, where a valid instrument would give one value. The size-menu instruments therefore do not identify the payment margin, and the difference estimator of this appendix, which would transport an instrument's information from the subsample to the full sample, has nothing valid to transport. The fixed-effects estimate of the payment margin stands on the institutional argument and the sign argument of Section \ref{sec:identification}, its agreement with the within-borrower estimate, and its agreement with the published payment responses at the same margin (Table \ref{tab:benchmark}).

\subsection{The balance held fixed: the rate margin}\label{appendix:ratemargin}

The paper's payment margin moves the balance with the payment at a fixed term. With $\log m$, $\log T$ and $\log L$ all in the regression, the amortization identity $m = L\,r/(1 - (1+r)^{-T})$ ties the three, and the only remaining source of variation in any one of them at fixed values of the other two is the contract rate $r$: both coefficients then measure the response of default to the rate, and within a borrower the rate is the lender's price of that loan's risk at that origination. On the multi-loan borrowers of the ten-percent archive the within-borrower payment elasticity with the log balance held fixed is $1.99$ at twenty-four months against $0.53$ with the balance moving, and the term elasticity $1.62$ against $0.30$; the ratio of the two, near $0.8$, sits above the formula's ceiling in every group and is a ratio of two rate responses rather than a far-to-near payment weight. The balance-fixed estimate is the bureau's counterpart to the published quasi-experiments that move a payment through the rate, and Table \ref{tab:benchmark} reports it beside them; the paper's own margin stays the payment with the balance moving, the cleanest payment variation the bureau offers within borrower.

\subsection{The estimate on the auto panel}\label{appendix:ivvalidation:auto}

Section \ref{sec:identification} states the assumption under which the fixed-effects estimate of the payment margin is causal and reports that the size-menu instruments give the opposite sign. This subsection reports the difference estimator applied to the auto panel. The instruments are constructed on the zip3-by-origination-year cells of the covariate-matched extract, for lender-years with at least fifty originations, which hold 97 percent of loans: the lender's leave-one-out share of long-term originations in the year, which shifts $T$ through the maturity menu \citep{argyle_monthly_2020, hertzberg_screening_2018}, and the lender's leave-one-out mean log loan size, which shifts $m$ given $T$ through the size menu. On this subsample $S$ the two-stage least-squares estimate of $(\gamma_m, \gamma_T)$ is compared with the fixed-effects estimate on the same loans, and equation \eqref{eq:diffapp} corrects the full-sample fixed-effects estimate coefficient by coefficient. The estimator is consistent for the full-sample coefficients under equi-confounding, stated for this application as follows.

\begin{quote}
\emph{Equi-confounding.} Within a zip3-by-origination-year cell, the projection of the unabsorbed determinants of default on $(\log m, \log T)$ is the same for loans at large lenders as for loans at small lenders. Default rates, the sizes of the payment and term effects, and the distribution of contracts may all differ between the two groups; what may not differ is the within-cell relationship between contract terms and the risk the cell does not absorb.
\end{quote}

\noindent This is the assumption under which the instrument's information travels from the subsample to the rest of the data \citep{li_identification_2024, angrist_leveraging_2017}, and it is weaker than the assumption that large-lender loans are a random draw \citep{hotz_predicting_2005}. What it excludes is a confounder that operates only at small lenders: a pricing practice that ties the payment to unobserved risk in a way that large lenders' underwriting does not. Its testable implication is that the estimated bias, $\hat\gamma^{FE}(S) - \hat\gamma^{IV}(S)$, is the same across groups of loans that exist within $S$. Table \ref{tab:identification} reports the estimated bias for the subsample as a whole and for the groups within it that hold enough instrument variation to estimate it separately, the bottom income quartile and the South and West census regions; the standard errors on every row come from a cluster bootstrap over lenders, stratified by subsample, so that the covariance among the three pieces of equation \eqref{eq:diffapp} enters the corrected estimate.

The instruments are leave-out means, the construction of the examiner design \citep{kling_incarceration_2006, dobbie_targeted_2020}, with one difference that fixes where the argument can fail. Cases are randomly assigned to examiners; borrowers are not randomly assigned to lenders. The exclusion restriction therefore requires that, within a cell and within a lender, the year-to-year movement of the lender's menu is unrelated to the unobserved risk of the borrowers it attracts. The failure to guard against is sorting: a lender that expands long maturities in a year in which it also reaches down the risk distribution. Three features of the design address it. Lender fixed effects remove the cross-lender association between menus and clienteles, so that only within-lender changes in the menu identify; Table \ref{tab:identification} reports the corrected estimate with and without them. The two instruments over-identify the two coefficients by one restriction; a violation common to both instruments escapes a test of that restriction, a violation specific to one does not. And a confounder that moves default through both margins in proportion, as a lender's overall reach down the risk distribution does, scales $\gamma_m$ and $\gamma_T$ together; only a confounder that acts on one margin and not the other moves their relative size, which is a narrower class than the class of confounders of either coefficient alone.

Table \ref{tab:identification} reports the result. The fixed-effects payment coefficient is $+0.122$ (standard error $0.044$, linear-probability units times 100) on all loans and $+0.107$ ($0.045$) on the subsample; the two-stage least-squares coefficient on the subsample is $-0.713$ ($0.293$), so the difference estimator returns $-0.698$ ($0.292$), and with lender effects $-0.912$ ($0.502$). The instrument-corrected payment coefficient has the opposite sign to the fixed-effects one and to every payment response in the paper, including the within-borrower estimate and the rule-based experiments. The reason is in the instruments. Lender-years with larger loans serve safer borrowers, so the size-menu instrument moves the payment together with unobserved risk and the exclusion restriction fails; the estimated bias in Panel B, which equi-confounding requires to be the same across groups, is $0.42$ in the South and $1.49$ in the West. The size-menu instruments therefore do not identify the payment margin (Appendix \ref{appendix:strategies}), and the paper's estimate is the fixed-effects one, defended as Section \ref{sec:identification} states. Two limits of the difference estimator remain in any application of it to these data. Groups of borrowers that exist only at small lenders receive a correction estimated at large lenders; for them equi-confounding is an assumption without a test. And the instruments identify the response of compliers to the lender's menu, whereas the fixed-effects coefficients weight every loan by its within-cell variance; the bias term in \eqref{eq:diffapp} absorbs the difference between these two averages along with any confounding, which is the cost of using the instrument as a level rather than as a test.

\begin{table}[ht]
\centering
\caption{The payment margin: fixed effects, instruments, and the corrected estimate on the large-lender subsample}
\label{tab:identification}
\resizebox{\textwidth}{!}{\begin{tabular}{lrrrr}
\toprule
Design & $\gamma_m$ & (s.e.) & $\gamma_T$ & (s.e.) \\
\midrule
\multicolumn{5}{l}{\emph{Panel A: default within twelve months (linear probability coefficients $\times 100$)}} \\
Fixed effects, all loans & 0.122 & (0.044) & -0.836 & (0.188) \\
Fixed effects, large-lender subsample & 0.107 & (0.045) & -0.794 & (0.206) \\
Two-stage least squares, subsample & -0.713 & (0.293) & -0.717 & (0.485) \\
Difference estimator & -0.698 & (0.292) & -0.759 & (0.474) \\
Fixed effects, all loans, with lender effects & 0.163 & (0.026) & -0.093 & (0.095) \\
Fixed effects, subsample, with lender effects & 0.169 & (0.026) & -0.070 & (0.101) \\
2SLS, subsample, with lender effects & -0.906 & (0.502) & 1.749 & (0.821) \\
Difference estimator, with lender effects & -0.912 & (0.502) & 1.725 & (0.819) \\
\midrule
\multicolumn{5}{l}{\emph{Panel B: the estimated bias $\hat\gamma^{FE}(S) - \hat\gamma^{IV}(S)$ by group ($\times 100$)}} \\
All loans in the subsample & 0.809 & (0.037) & -0.052 & (0.045) \\
Score decile 1 & 1.021 & (0.242) & -0.528 & (0.214) \\
Score decile 2 & 0.426 & (0.149) & 0.114 & (0.159) \\
Score decile 3 & 0.512 & (0.094) & -0.085 & (0.109) \\
Score decile 4 & 0.294 & (0.063) & 0.142 & (0.092) \\
Score decile 5 & 0.250 & (0.044) & -0.140 & (0.058) \\
Score decile 6 & 0.075 & (0.024) & -0.067 & (0.039) \\
Score decile 7 & 0.065 & (0.023) & -0.042 & (0.030) \\
Income quartile 1 & 1.147 & (0.125) & 0.043 & (0.114) \\
Income quartile 2 & 0.686 & (0.055) & -0.640 & (0.059) \\
Income quartile 3 & 0.171 & (0.030) & -0.244 & (0.033) \\
Income quartile 4 & 0.059 & (0.015) & -0.056 & (0.020) \\
Region NE & 0.635 & (0.061) & -0.324 & (0.062) \\
Region MW & 0.248 & (0.054) & -0.074 & (0.085) \\
Region S & 0.888 & (0.076) & -0.001 & (0.085) \\
Region W & 1.188 & (0.064) & 0.209 & (0.095) \\
\bottomrule
\end{tabular}
}
\vspace{0.5em}
\begin{minipage}{0.95\textwidth}
\small \textit{Notes:} Rows report the coefficients on $\log m$ and $\log T$ for default within twelve months on the covariate-matched extract: fixed effects on all of its loans, fixed effects and two-stage least squares on the subsample of lender-years with at least fifty originations, the estimated bias, and the difference estimator of equation \eqref{eq:diffapp}, with and without lender fixed effects; standard errors from a cluster bootstrap over lenders stratified by subsample. The lower panel reports the estimated bias for the subsample and for the groups within it with enough instrument variation to estimate it separately.
\end{minipage}
% replication: Code/identification/I2c_margins_iv.R, I2d_difference_estimator.R
\end{table}

\paragraph{The within-borrower specification with lender-year cells.} Every within-borrower estimate in the paper absorbs borrower and origination-year effects, so that a borrower's contracts are compared across lenders at the same date. The same regression on the same loans with lender-year in place of origination-year effects, which holds the lender's menu fixed as well as the borrower, gives a payment elasticity of $0.51$ ($0.005$) within twelve months, $0.57$ ($0.003$) within twenty-four and $0.47$ over the life of the loan, and a term elasticity of $-0.19$ ($0.01$), $0.00$ ($0.01$) and $0.42$ ($0.01$): the twelve-month term response and the lifetime response survive, and the twenty-four-month response, which the origination-year specification puts at $0.30$, is absent. Within one lender's menu in one year the contract length varies with the amount financed and the borrower's screening by that lender \citep{hertzberg_screening_2018}, so the two specifications compare different pairs of contracts, and the paper reports both. Replication: \texttt{Code/identification/I6\_within\_person.R}, rows \texttt{borrower + lender-year} of \texttt{results/identification\_within\_person.csv}.

